\PassOptionsToPackage{expansion=false}{microtype} % disable font expansion to avoid bitmap font errors
\documentclass[journal]{new-aiaa}
\ifdefined\pdfmapfile
  \pdfmapfile{+qtm.map} % load TeX Gyre Termes map (not included by default in this MiKTeX setup)
  \pdfmapfile{+newtx.map} % ensure newtx fonts use bundled map data
\fi
\usepackage{graphicx}
\graphicspath{ {./figures/} }
\usepackage[rightcaption]{sidecap}
\usepackage{wrapfig}
\usepackage{float}
\usepackage{imakeidx}
\usepackage{comment}
\usepackage{commath}
\usepackage[titletoc]{appendix}
\usepackage{graphicx}
\usepackage{resizegather}
\usepackage[english]{babel}
\usepackage{tabu}
\usepackage{booktabs}
\usepackage{xfrac}
\usepackage{tabularx}
\usepackage{amsmath}
\usepackage{commath}
\usepackage{graphicx,bm}
\usepackage{verbatim}
\usepackage{lscape}
\usepackage{relsize}
\usepackage{enumitem}
\usepackage{textcomp}
\usepackage{breqn}
\usepackage{makecell}
\usepackage{longtable,tabularx}
\usepackage{multirow}
\usepackage{doi}
\usepackage{fancyhdr}
\usepackage{algorithm}
\usepackage{algpseudocode}
\usepackage{setspace}
\usepackage{footnote}
\PassOptionsToPackage{hyphens}{url}
\usepackage{hyperref}
\hypersetup{colorlinks,linkcolor={blue},citecolor={blue},urlcolor={blue}}
\usepackage[numbers]{natbib}
\usepackage{mathtools}
\usepackage[disable]{todonotes}
\usepackage[framemethod=tikz]{mdframed}
\usepackage{booktabs,xcolor,siunitx}
\usepackage{soul}
\usepackage{cleveref}
\renewcommand{\ttdefault}{lmtt} % ensure scalable typewriter font for \texttt

% Custom commands
\newcommand{\ra}[1]{\renewcommand{\arraystretch}{#1}}
\newcommand{\degree}{\ensuremath{^\circ\,}}
\newcommand{\overbar}[1]{\mkern 1.5mu\overline{\mkern-1.5mu#1\mkern-1.5mu}\mkern 1.5mu}
\newcommand{\f}[2]{\frac{#1}{#2}}
\newcommand{\mb}[1]{\mathbf{#1}}
\newcommand{\tr}[1]{\mathrm{Tr}\left({#1}\right)}
\newcommand{\mbg}[1]{\bm{#1}}
\DeclarePairedDelimiter{\ceil}{\lceil}{\rceil}
\DeclareMathAlphabet\mathbfcal{OMS}{cmsy}{b}{n}
\renewcommand{\d}{\mathop{}\!\mathrm{d}} % total derivative
\newcommand{\p}{\partial}

\title{Scalable Time-Accurate Adjoint Sensitivity Analysis \\for Unsteady Compressible Flows}

\author{Sicheng He\footnote{Assistant Professor, AIAA member, sicheng@utk.edu}}
\affil{University of Tennessee, Department of Mechanical and Aerospace Engineering, Knoxville, TN 37996}

\date{\today}

\begin{document}
\maketitle{}


\begin{abstract}
Gradient-based optimization of unsteady aerodynamic flows requires efficient sensitivity analysis, yet the cost of finite-difference approaches scales linearly with the number of design variables, making them impractical for high-dimensional problems.
We present a discrete adjoint framework for time-accurate compressible flow simulations that leverages PETSc's time-stepping adjoint infrastructure to compute sensitivities at a cost independent of the number of design variables.
The framework uses the Crank--Nicolson scheme for time integration, matrix-free Jacobian-transpose--vector products from reverse-mode algorithmic differentiation, and revolve checkpointing to manage memory for large-scale problems.
We implement this approach in ADflow, an open-source finite-volume Reynolds-averaged Navier--Stokes solver, and support both final-time and time-integrated objective functions.
The adjoint sensitivities are verified against finite differences on three test cases of increasing complexity: vortex shedding behind a circular cylinder, a pitching airfoil, and transonic buffet over a supercritical wing section.
% placeholder: add quantitative accuracy statement
We demonstrate the practical utility of the framework by performing unsteady aerodynamic shape optimization.
% placeholder: add optimization result summary
\end{abstract}

\section{Introduction}\label{sec:introduction}
Adjoint-based sensitivity analysis has become an indispensable tool for aerodynamic shape optimization~\cite{Pironneau1974, Jameson1988, Martins2013}.
For steady flows, the adjoint method computes the gradient of an objective function with respect to all design variables at a cost comparable to a single flow solve, regardless of the number of design variables.
This favorable scaling has enabled large-scale aerodynamic and multidisciplinary design optimization with hundreds or thousands of shape parameters~\cite{Kenway2019, He2018}.

Extending adjoint methods to unsteady flows introduces additional challenges.
The forward flow solution must be available at every time step during the backward adjoint integration, which creates substantial memory requirements for time-accurate simulations with many time steps.
Furthermore, the adjoint equations form a system that must be marched backward in time, requiring transposed Jacobian operations at each step.
These challenges have motivated a substantial body of work on unsteady adjoint methods, spanning both continuous~\cite{Nadarajah2003, Rumpfkeil2010} and discrete~\cite{Nielsen2004, Vishnampet2015} formulations.

The continuous adjoint approach derives the adjoint equations at the partial differential equation (PDE) level and then discretizes them.
While this avoids differentiating through the discretization, the resulting adjoint solution may not be consistent with the primal discretization, leading to inaccurate gradients unless dual consistency is carefully maintained~\cite{Hicken2014}.
The discrete adjoint approach differentiates the fully discretized system, guaranteeing gradient accuracy to machine precision.
Algorithmic differentiation (AD) tools such as Tapenade~\cite{Hascoet2013} can generate the required derivative code systematically, making the discrete approach attractive for complex solvers.

Several unsteady adjoint implementations have been reported for compressible flow solvers.
Nadarajah and Jameson~\cite{Nadarajah2003} developed both continuous and discrete adjoint methods for the Euler equations with a dual time-stepping scheme.
Rumpfkeil and Zingg~\cite{Rumpfkeil2010} presented a general framework for optimal control of unsteady flows using the discrete adjoint, with applications to the Euler and Navier--Stokes equations.
Nielsen and Diskin~\cite{Nielsen2004} demonstrated a discrete adjoint for unsteady turbulent flows on dynamic unstructured grids.
Vishnampet et al.~\cite{Vishnampet2015} developed a practical discrete adjoint for high-fidelity compressible turbulence simulations.
Despite these advances, most existing implementations are tightly coupled to specific solvers and time-integration schemes, limiting code reuse and extensibility.

PETSc (Portable, Extensible Toolkit for Scientific Computation)~\cite{petsc-web-page, petsc-user-ref} provides a mature infrastructure for solving ordinary differential equations (ODEs) and differential-algebraic equations (DAEs) through its time-stepping (TS) component.
Zhang et al.~\cite{Zhang2022} extended PETSc's TS module with a discrete adjoint capability that supports first- and second-order sensitivity analysis for the implicit theta family of time integrators.
This infrastructure handles the adjoint equation derivation, backward time marching, and trajectory management internally, allowing solver developers to obtain adjoint sensitivities by providing only the residual function and its Jacobian.
The revolve algorithm~\cite{Griewank2000} is integrated for optimal checkpointing, and the petsc4py interface~\cite{Dalcin2011} enables access from Python.

In this work, we leverage PETSc's TS adjoint infrastructure to develop a time-accurate discrete adjoint capability for ADflow~\cite{ADflow2020}, an open-source structured-multiblock finite-volume solver for the compressible RANS equations.
ADflow already provides a well-validated steady-state adjoint solver~\cite{Kenway2019, Mader2020} based on AD-generated Jacobian products and PETSc's Krylov solvers.
We extend this to the unsteady regime by casting ADflow's spatial residual evaluation as PETSc TS callbacks, using the Crank--Nicolson scheme for second-order temporal accuracy, and employing matrix-free operations via PETSc's \texttt{MatShell} to avoid explicit Jacobian assembly.
Both final-time and time-integrated objective functions are supported through PETSc's quadrature time-stepping interface.
We verify the adjoint sensitivities against finite differences on three test cases---cylinder vortex shedding, a pitching airfoil, and transonic buffet---and demonstrate the framework in an unsteady aerodynamic shape optimization.

The remainder of this paper is organized as follows.
\Cref{sec:governing} presents the governing equations and their spatial discretization.
\Cref{sec:adjoint} describes the discrete adjoint formulation for time-accurate problems, including the time-integrated objective, matrix-free implementation, and checkpointing.
\Cref{sec:implementation} details the implementation in ADflow.
\Cref{sec:verification} presents verification results, and \Cref{sec:demonstration} demonstrates unsteady shape optimization.
\Cref{sec:conclusions} provides concluding remarks.

\section{Governing Equations and Spatial Discretization}\label{sec:governing}
The compressible RANS equations express the conservation of mass, momentum, and energy for a fluid element.
In integral form over a control volume $\Omega$ with boundary $\p \Omega$, the equations read
\begin{equation}\label{eq:rans_integral}
    \f{\p}{\p t} \int_{\Omega} \mb{q} \, \d V + \oint_{\p \Omega} \left( \mb{f}_c - \mb{f}_v \right) \cdot \hat{\mb{n}} \, \d S = \mb{0},
\end{equation}
where $\mb{q} = [\rho, \, \rho u, \, \rho v, \, \rho w, \, \rho E]^\intercal \in \mathbb{R}^5$ is the vector of conservative variables, $\mb{f}_c$ and $\mb{f}_v$ are the inviscid and viscous flux vectors, and $\hat{\mb{n}}$ is the outward unit normal to $\p \Omega$.
Here $\rho$ is the density, $(u, v, w)$ are the Cartesian velocity components, and $E$ is the total energy per unit mass.
The system is closed with the ideal gas equation of state, $p = \rho (\gamma - 1)(E - (u^2 + v^2 + w^2)/2)$, where $p$ is the pressure and $\gamma$ is the ratio of specific heats.

Turbulence is modeled using the SA one-equation model~\cite{Spalart1992}, which introduces an additional transport equation for the modified turbulence viscosity $\tilde{\nu}$.
The SA equation is solved in a fully coupled manner with the mean flow equations, resulting in a system of six conservation laws per cell ($m = 6$).
The eddy viscosity $\mu_t$ is computed from $\tilde{\nu}$ and enters the viscous fluxes through the Reynolds stress tensor (via the Boussinesq approximation) and the turbulent heat flux.

ADflow~\cite{ADflow2020} discretizes \cref{eq:rans_integral} on structured multiblock grids using a cell-centered finite-volume method.
The inviscid fluxes are computed using central differences with added scalar artificial dissipation for stability.
The viscous fluxes are discretized using a second-order central scheme.
The overall spatial discretization is second-order accurate.
After spatial discretization, the governing equations reduce to a system of ODEs in semi-discrete form,
\begin{equation}\label{eq:semi_discrete}
    \f{\d \mb{w}}{\d t} = \mb{r}(\mb{w}, \mb{x}),
\end{equation}
where $\mb{w} \in \mathbb{R}^n$ is the vector of conservative flow variables at all $n/m$ cells (with $m = 6$ variables per cell for RANS), $\mb{x} \in \mathbb{R}^{n_x}$ denotes the design parameters (e.g., shape variables or flow conditions), and $\mb{r}: \mathbb{R}^n \times \mathbb{R}^{n_x} \to \mathbb{R}^n$ is the spatial residual operator that encapsulates the inviscid fluxes, viscous fluxes, and turbulence source terms.

For problems with moving or deforming grids (e.g., pitching airfoils), the governing equations are formulated in an arbitrary Lagrangian--Eulerian (ALE) frame.
The grid velocities enter the convective fluxes, and the geometric conservation law (GCL) is satisfied to preserve freestream preservation on moving meshes.
In this case, the residual depends additionally on the mesh coordinates and their time derivatives, both of which are functions of the design parameters $\mb{x}$.

\section{Discrete Adjoint for Time-Accurate Problems}\label{sec:adjoint}
This section describes the discrete adjoint formulation for time-accurate simulations, including the forward time integration scheme, the adjoint derivation for both final-time and time-integrated objectives, the matrix-free implementation, and the checkpointing strategy.

\subsection{Forward Time Integration}\label{sec:forward}
We cast the semi-discrete system~\eqref{eq:semi_discrete} in implicit form as required by PETSc's TS framework~\cite{Zhang2022},
\begin{equation}\label{eq:implicit_form}
    \mb{h}(t, \mb{w}, \dot{\mb{w}}) \coloneqq \dot{\mb{w}} - \mb{r}(\mb{w}, \mb{x}) = \mb{0},
\end{equation}
where $\dot{\mb{w}} \coloneqq \d \mb{w} / \d t$ is the time derivative of the state vector.
In PETSc's notation, the user provides the implicit residual $\mb{h}$ through an \texttt{IFunction} callback and the associated Jacobian $\p \mb{h}/\p \mb{w} + a \, \p \mb{h}/\p \dot{\mb{w}}$ through an \texttt{IJacobian} callback, where $a$ is a shift parameter determined by the time integration scheme.
For \cref{eq:implicit_form}, the Jacobian simplifies to $a \, \mb{I} - \p \mb{r}/\p \mb{w}$, where $\mb{I} \in \mathbb{R}^{n \times n}$ is the identity matrix.

We use the implicit theta method, which advances the solution from $t^n$ to $t^{n+1} = t^n + \Delta t$ as
\begin{equation}\label{eq:theta_method}
    \mb{w}^{n+1} = \mb{w}^n + \Delta t \left[\theta \, \mb{r}(\mb{w}^{n+1}, \mb{x}) + (1-\theta) \, \mb{r}(\mb{w}^n, \mb{x})\right],
\end{equation}
where $\Delta t$ is the time step size and $\theta \in [0, 1]$ controls the degree of implicitness.
Setting $\theta = 1$ recovers the first-order backward Euler method, while $\theta = 1/2$ yields the second-order Crank--Nicolson (CN) scheme.
We use CN ($\theta = 1/2$) throughout this work for its second-order temporal accuracy and A-stability.
At each time step, the nonlinear system arising from \cref{eq:theta_method} is solved using a Newton--Krylov method.
Defining the nonlinear residual for the $n$-th time step as $\mb{g}(\mb{w}^{n+1}) \coloneqq \mb{w}^{n+1} - \mb{w}^n - \Delta t [\theta \, \mb{r}(\mb{w}^{n+1}, \mb{x}) + (1-\theta) \, \mb{r}(\mb{w}^n, \mb{x})]$, the Newton update $\delta \mb{w}$ at each iteration is obtained by solving the linear system
\begin{equation}\label{eq:newton_system}
    \left(\mb{I} - \theta \, \Delta t \, \f{\p \mb{r}}{\p \mb{w}}\right) \delta \mb{w} = -\mb{g}(\mb{w}^{n+1}),
\end{equation}
where $\mb{I} - \theta \, \Delta t \, \p \mb{r}/\p \mb{w}$ is the Jacobian of the nonlinear system.
This linear system is solved iteratively using a preconditioned Krylov method (GMRES).
ADflow constructs a preconditioner from an approximate Jacobian assembled using first-order spatial discretization with scalar dissipation, which is factored with an incomplete LU (ILU) decomposition within an additive Schwarz method (ASM) framework.
The matrix-vector products for GMRES are computed exactly using forward-mode AD (Tapenade), so only the preconditioner uses the approximate Jacobian.
The mean flow and turbulence equations are solved in a fully coupled manner; that is, the state vector $\mb{w}$ and the Jacobian include both the mean flow and the SA turbulence variable at every Newton iteration.

\subsection{Adjoint Time Integration}\label{sec:adjoint_time}
We consider two types of objective functions.
The first depends only on the solution at the final time $t_f$,
\begin{equation}\label{eq:objective_final}
    f_{\text{final}} = \phi(\mb{w}(t_f), \mb{x}),
\end{equation}
where $\phi: \mathbb{R}^n \times \mathbb{R}^{n_x} \to \mathbb{R}$ is a scalar function (e.g., the lift or drag coefficient at the final time).
The second is a time-integrated objective,
\begin{equation}\label{eq:objective_integral}
    f_{\text{int}} = \int_{t_0}^{t_f} r(\mb{w}(t), \mb{x}, t) \, \d t,
\end{equation}
where $r: \mathbb{R}^n \times \mathbb{R}^{n_x} \times \mathbb{R} \to \mathbb{R}$ is an instantaneous cost rate (e.g., the time-averaged drag coefficient is obtained by setting $r = C_D(\mb{w}, \mb{x}) / (t_f - t_0)$).
The general objective combines both contributions,
\begin{equation}\label{eq:objective_general}
    f = \phi(\mb{w}(t_f), \mb{x}) + \int_{t_0}^{t_f} r(\mb{w}(t), \mb{x}, t) \, \d t.
\end{equation}

The discrete adjoint computes the total sensitivity $\d f / \d \mb{x}$ by integrating the adjoint equations backward in time.
PETSc's TS adjoint framework~\cite{Zhang2022} derives the adjoint equations by differentiating the discrete time-stepping scheme.
The adjoint variable $\mbg{\lambda}^n \in \mathbb{R}^n$ is initialized at the final time with
\begin{equation}\label{eq:adjoint_terminal}
    \mbg{\lambda}^{n_t} = \left(\f{\p \phi}{\p \mb{w}}\bigg|^{n_t}\right)^\intercal,
\end{equation}
where $n_t$ is the total number of time steps, and then marched backward.
For the theta method, the adjoint update at step $n$ requires solving the transposed linear system
\begin{equation}\label{eq:adjoint_system}
    \left(\mb{I} - \theta \, \Delta t \, \f{\p \mb{r}}{\p \mb{w}}\bigg|^{n}\right)^\intercal \mbg{\lambda}^{n}_{\text{stage}} = \mbg{\lambda}^{n+1} + \Delta t \left(\f{\p r}{\p \mb{w}}\bigg|^n\right)^\intercal,
\end{equation}
where $\mbg{\lambda}^{n}_{\text{stage}}$ is the stage adjoint variable and the second term on the right-hand side is present only when time-integrated objectives are used.
The adjoint variable is then updated as
\begin{equation}\label{eq:adjoint_update}
    \mbg{\lambda}^{n} = \mbg{\lambda}^{n}_{\text{stage}} + (1-\theta) \, \Delta t \left(\f{\p \mb{r}}{\p \mb{w}}\bigg|^{n}\right)^\intercal \mbg{\lambda}^{n}_{\text{stage}}.
\end{equation}
The parameter sensitivity is accumulated over all time steps as
\begin{equation}\label{eq:param_sens}
    \f{\d f}{\d \mb{x}} = \f{\p \phi}{\p \mb{x}}\bigg|^{n_t} + \sum_{n=0}^{n_t} \Delta t \left[\theta \left(\f{\p \mb{r}}{\p \mb{x}}\bigg|^{n+1}\right)^\intercal + (1-\theta) \left(\f{\p \mb{r}}{\p \mb{x}}\bigg|^{n}\right)^\intercal\right] \mbg{\lambda}^{n}_{\text{stage}} + \sum_{n=0}^{n_t} \Delta t \, \f{\p r}{\p \mb{x}}\bigg|^{n}.
\end{equation}
PETSc handles the derivation, backward time marching, and accumulation internally; the user provides only the forward residual $\mb{h}$, the Jacobians $\p \mb{h}/\p \mb{w}$ and $\p \mb{h}/\p \mb{x}$, and the objective function derivatives.

The time-integrated objective~\eqref{eq:objective_integral} is evaluated during the forward solve using a quadrature sub-integrator created by PETSc's \texttt{TSCreateQuadratureTS} interface.
This sub-integrator advances a separate set of variables that accumulate the integral, and its adjoint is handled automatically during the backward sweep.

The adjoint equations~\eqref{eq:adjoint_system}--\eqref{eq:param_sens} are derived under the assumption that the nonlinear system at each forward time step is solved exactly, i.e., $\mb{g}(\mb{w}^{n+1}) = \mb{0}$.
In practice, the Newton--Krylov solver converges the residual to a finite tolerance $\|\mb{g}\| \leq \epsilon$.
If the forward residual is $O(\epsilon)$, the adjoint gradient inherits an $O(\epsilon)$ error because the adjoint equations correspond to the exact discrete map rather than the approximate one realized by the inexact solve.
This error can accumulate over the $n_t$ time steps, so the forward solve tolerance must be chosen sufficiently tight to ensure that the adjoint sensitivities agree with finite differences to the desired accuracy.
We investigate this trade-off numerically in \cref{sec:verification}.

\subsection{Matrix-Free Implementation}\label{sec:matfree}
The transposed linear system~\eqref{eq:adjoint_system} requires the Jacobian-transpose--vector product $(\p \mb{r}/\p \mb{w})^\intercal \mbg{\lambda}$ at each adjoint time step.
For large-scale problems, explicit assembly and storage of the $n \times n$ Jacobian matrix is prohibitively expensive.
ADflow provides matrix-free Jacobian-transpose--vector products generated by applying the reverse mode of AD (via Tapenade~\cite{Hascoet2013}) to the residual evaluation routine.
These products are exposed to PETSc through a \texttt{MatShell} object, which implements the \texttt{MATOP\_MULT\_TRANSPOSE} operation.

When the \texttt{IJacobian} callback is invoked with shift parameter $a$, the \texttt{MatShell} stores the current state $\mb{w}^n$ and the shift $a$.
Subsequent calls to \texttt{MatMultTranspose} compute
\begin{equation}\label{eq:matshell}
    \mb{y} = \left(a \, \mb{I} - \f{\p \mb{r}}{\p \mb{w}}\right)^\intercal \mb{v} = a \, \mb{v} - \left(\f{\p \mb{r}}{\p \mb{w}}\right)^\intercal \mb{v},
\end{equation}
where the second term is evaluated by ADflow's reverse-mode AD routine using the stored state $\mb{w}^n$, and $\mb{v} \in \mathbb{R}^n$ is the input vector.

A preconditioner is required for efficient convergence of the Krylov solver within each adjoint time step.
ADflow provides an approximate Jacobian assembled using first-order spatial discretization with scalar dissipation, which serves as the preconditioner matrix.
This approximate Jacobian is passed as the second argument to PETSc's \texttt{TSSetIJacobian}, allowing PETSc to use it for preconditioning while the matrix-free \texttt{MatShell} is used for the actual matrix-vector products.
The approximate Jacobian is factored using an incomplete LU (ILU) decomposition within an additive Schwarz method (ASM) domain decomposition framework.

Similarly, the parameter Jacobian-transpose--vector product $(\p \mb{r}/\p \mb{x})^\intercal \mbg{\lambda}$ needed in \cref{eq:param_sens} is computed using ADflow's reverse-mode AD with respect to the design variables.
This product is provided to PETSc through the \texttt{IJacobianP} callback.

\subsection{Checkpointing}\label{sec:checkpointing}
The adjoint time integration requires access to the forward solution $\mb{w}^n$ at every time step $n = n_t, n_t-1, \ldots, 0$.
For large-scale three-dimensional RANS simulations, each state snapshot can require hundreds of megabytes, making storage of the full trajectory infeasible for long time horizons.
For example, a simulation with $5 \times 10^6$ cells, six variables per cell, and 1000 time steps requires approximately 240~GB of storage.

We address this through the revolve checkpointing algorithm~\cite{Griewank2000}, which is integrated into PETSc's \texttt{TSTrajectory} interface.
Revolve stores a limited number of checkpoints in memory and recomputes intermediate states as needed during the backward sweep.
For $n_t$ time steps and $c$ available checkpoint slots, the optimal schedule results in an additional computational cost proportional to $\log(n_t / c)$ forward solves.
The number of checkpoint slots is specified at runtime, allowing the user to balance memory usage against recomputation cost.

\section{Implementation in ADflow}\label{sec:implementation}
ADflow~\cite{ADflow2020} is a structured-multiblock finite-volume solver for the compressible RANS equations.
It is written primarily in Fortran 90 with a Python interface provided through f2py, and it uses PETSc for linear algebra operations.
ADflow has a well-established steady-state adjoint solver~\cite{Kenway2019} that uses AD-generated (Tapenade) Jacobian products and PETSc Krylov solvers.
This section describes how we extend ADflow to support time-accurate adjoint sensitivity analysis by interfacing with PETSc's TS module through petsc4py~\cite{Dalcin2011}.

The time-accurate adjoint integration is driven from the Python level.
The key implementation steps are as follows.

First, a PETSc TS object is created and configured with the Crank--Nicolson scheme (\texttt{TSCN}).
The time step size, final time, and convergence tolerances are set based on the user's input parameters.
Trajectory saving is enabled via \texttt{TSSetSaveTrajectory} to store checkpoints for the adjoint solve, and the revolve schedule is configured through command-line options.

Second, the \texttt{IFunction} callback wraps ADflow's spatial residual evaluation.
Given the current state $\mb{w}$ and time derivative $\dot{\mb{w}}$, it evaluates $\mb{h} = \dot{\mb{w}} - \mb{r}(\mb{w}, \mb{x})$ by calling ADflow's Fortran-level residual routine and subtracting the result from $\dot{\mb{w}}$.

Third, the \texttt{IJacobian} callback creates a \texttt{MatShell} for the matrix-free Jacobian $a \, \mb{I} - \p \mb{r}/\p \mb{w}$ and assembles ADflow's approximate Jacobian for preconditioning.
The \texttt{MatShell}'s \texttt{MultTranspose} operation calls ADflow's Tapenade-generated reverse-mode routine to evaluate $(\p \mb{r}/\p \mb{w})^\intercal \mb{v}$ and adds the shift contribution $a \, \mb{v}$.

Fourth, for parameter sensitivities, the \texttt{IJacobianP} callback provides $(\p \mb{h}/\p \mb{x})^\intercal \mbg{\lambda}$ using ADflow's reverse-mode AD with respect to the design variables.
This includes sensitivities with respect to mesh coordinates (for shape optimization) and flow conditions (e.g., angle of attack, Mach number).

Fifth, for time-integrated objectives, a quadrature sub-integrator is created via \texttt{TSCreateQuadratureTS}.
The cost rate function $r(\mb{w}, \mb{x}, t)$ and its derivatives with respect to $\mb{w}$ and $\mb{x}$ are provided through the sub-integrator's callbacks.

The overall workflow proceeds as follows: (1) compute the initial condition (e.g., a converged steady-state solution), (2) perform the forward solve via \texttt{TSSolve}, (3) set the adjoint terminal conditions via \texttt{TSSetCostGradients}, (4) execute the backward adjoint solve via \texttt{TSAdjointSolve}, and (5) extract the accumulated parameter sensitivities.
\Cref{alg:workflow} summarizes this procedure.

\begin{algorithm}[t]
\caption{Time-accurate adjoint workflow in ADflow.}\label{alg:workflow}
\begin{algorithmic}[1]
\State $\mb{w}_0 \gets \text{SteadySolve}()$ \Comment{base flow}
\State $\text{TS} \gets \text{CreateTS}(\theta = 1/2)$ \Comment{Crank--Nicolson}
\State $\text{RegisterCallbacks}(\texttt{IFunction}, \texttt{IJacobian}, \texttt{IJacobianP})$
\State $\text{EnableCheckpointing}(\text{revolve})$
\State $\mb{w}(t_f) \gets \texttt{TSSolve}(\mb{w}_0)$ \Comment{forward solve}
\State $\mbg{\lambda}^{n_t} \gets (\p \phi / \p \mb{w})^\intercal$ \Comment{terminal condition}
\State $\mbg{\mu} \gets (\p \phi / \p \mb{x})^\intercal$ \Comment{parameter seed}
\State $\mbg{\lambda}^0, \mbg{\mu} \gets \texttt{TSAdjointSolve}()$ \Comment{backward sweep}
\State $\d f/\d \mb{x} = \mbg{\mu} + \mbg{\lambda}^0 \cdot (\d \mb{w}_0 / \d \mb{x})$ \Comment{total sensitivity}
\end{algorithmic}
\end{algorithm}

\section{Verification and Validation}\label{sec:verification}
We verify the adjoint sensitivities by comparing them against finite-difference (FD) approximations for three test cases of increasing complexity.
For each case, the relative error between the adjoint and FD sensitivities is computed as
\begin{equation}\label{eq:rel_error}
    \epsilon_{\text{rel}} = \f{\left| g_{\text{adj}} - g_{\text{FD}} \right|}{\left| g_{\text{FD}} \right|},
\end{equation}
where $g_{\text{adj}}$ and $g_{\text{FD}}$ denote the adjoint and FD sensitivity values, respectively.
The FD step size is selected to balance truncation and round-off errors; we use centered differences with a step size of $h = 10^{-5}$ unless otherwise noted.

\subsection{Cylinder Vortex Shedding}\label{sec:cylinder}
The first test case is the two-dimensional flow over a circular cylinder at a Reynolds number of $\text{Re} = 100$ and a Mach number of $M = 0.2$.
At this Reynolds number, the flow develops a periodic von K\'arm\'an vortex street downstream of the cylinder.
This case is solved using the laminar Navier--Stokes equations (no turbulence model).

The computational domain extends [placeholder: domain size] diameters in each direction from the cylinder center, and the grid contains [placeholder: cell count] cells.
The simulation is initialized from a steady-state solution and run for [placeholder: number] time steps with a time step size of $\Delta t = $ [placeholder: value].
The objective function is the time-averaged drag coefficient,
\begin{equation}\label{eq:cylinder_obj}
    f = \f{1}{t_f - t_0} \int_{t_0}^{t_f} C_D(\mb{w}(t)) \, \d t.
\end{equation}

% placeholder: sensitivity results table comparing adjoint vs FD
% placeholder: figure showing drag history and adjoint convergence

\subsection{Pitching Airfoil}\label{sec:pitching}
The second test case is a NACA 0012 airfoil undergoing prescribed pitching motion about its quarter-chord.
The freestream conditions are $M = 0.3$ and $\text{Re} = 1 \times 10^6$, and the pitching motion is
\begin{equation}\label{eq:pitch_motion}
    \alpha(t) = \alpha_0 + \alpha_1 \sin(2 \pi f \, t),
\end{equation}
where $\alpha_0$ is the mean angle of attack, $\alpha_1$ is the pitching amplitude, and $f$ is the pitching frequency.
The specific values are $\alpha_0 = 0\degree$, $\alpha_1 = 1\degree$, and $f =$ [placeholder: value]~Hz.
This case uses the SA turbulence model and an ALE formulation with mesh deformation to accommodate the airfoil motion.

The grid contains [placeholder: cell count] cells, and the simulation is run for [placeholder: number] periods with [placeholder: steps per period] time steps per period.
The objective function is the peak lift coefficient over the simulation window,
\begin{equation}\label{eq:pitch_obj}
    f = C_L(\mb{w}(t^*)),
\end{equation}
where $t^*$ is the time at which $C_L$ reaches its maximum.
Sensitivities are computed with respect to the airfoil shape parameters and the mean angle of attack $\alpha_0$.

% placeholder: sensitivity results table comparing adjoint vs FD
% placeholder: figure showing Cl history and adjoint-vs-FD comparison

\subsection{Transonic Buffet}\label{sec:buffet}
The third test case is transonic buffet over an OAT15A supercritical airfoil at $M = 0.73$ and $\alpha = 3.5\degree$.
At these conditions, a self-sustained shock oscillation develops on the upper surface, producing large fluctuations in the aerodynamic loads~\cite{Crouch2009}.
This case uses the SA turbulence model.

The grid contains [placeholder: cell count] cells, and the simulation is run for [placeholder: number] buffet cycles.
The time step size is chosen to resolve the buffet frequency with at least [placeholder: number] steps per cycle.
The objective function is the time-averaged drag coefficient, analogous to \cref{eq:cylinder_obj}.
Sensitivities are computed with respect to the airfoil shape parameters.

% placeholder: sensitivity results table comparing adjoint vs FD
% placeholder: figure showing buffet Cd history and shock position

\section{Demonstration: Unsteady Shape Optimization}\label{sec:demonstration}
To demonstrate the practical utility of the time-accurate adjoint framework, we perform an unsteady aerodynamic shape optimization.

% placeholder: describe optimization problem setup
% - which test case (pitching airfoil or buffet)
% - design variables (FFD control points)
% - objective function
% - constraints
% - optimizer (SNOPT or IPOPT via pyOptSparse)

% placeholder: optimization results
% - convergence history
% - shape comparison (initial vs optimized)
% - objective improvement
% - flow field comparison

\section{Conclusions}\label{sec:conclusions}
We developed a scalable time-accurate discrete adjoint framework for unsteady compressible flow simulations by interfacing ADflow with PETSc's time-stepping adjoint infrastructure.
The framework uses the Crank--Nicolson scheme for second-order temporal accuracy, matrix-free Jacobian-transpose--vector products from reverse-mode algorithmic differentiation via Tapenade, and revolve checkpointing to manage memory for large-scale problems.
Both final-time and time-integrated objective functions are supported through PETSc's quadrature time-stepping interface.

The adjoint sensitivities were verified against finite differences on three test cases: vortex shedding behind a circular cylinder, a pitching airfoil with mesh deformation, and transonic buffet over a supercritical airfoil.
% placeholder: summarize accuracy (e.g., relative errors below X for all cases)
The framework was demonstrated in an unsteady aerodynamic shape optimization.
% placeholder: summarize optimization result

By leveraging PETSc's TS adjoint module, the implementation avoids manually deriving and coding the adjoint time-integration scheme, reducing development effort and the risk of implementation errors.
The matrix-free approach enables application to large-scale three-dimensional problems without explicit Jacobian storage.
The revolve checkpointing strategy provides a practical balance between memory and computational cost, with the number of checkpoints configurable at runtime.

% placeholder: future work (e.g., second-order adjoint for Hessian, aeroelastic coupling, higher-order time schemes)

\bibliography{bib/references.bib,bib/agi.bib}

\end{document}
